% Formation du diiode : concentration en fonction du temps (exercice 2, chapitre 4)
\begin{tikzpicture}
\begin{axis}[otMesures,otAxesBout,otLecture,width=104mm,height=46mm,xmin=0,xmax=92,ymin=0,ymax=5,
  xtick={0,10,...,90},ytick={0,1,...,5},xlabel={$t$ (min)},ylabel={$[\mathrm{I_2}]$ (mmol$\cdot$L$^{-1}$)}]
\addplot[otNavy,line width=1pt,domain=0:90,samples=120] {4*(1-2^(-x/15))};
\solution{\draw[otMuted,line width=0.5pt,dash pattern=on 1.5pt off 1pt] (axis cs:0,4) -- (axis cs:90,4);
\draw[otGreenDark,line width=0.6pt,dash pattern=on 2pt off 1.5pt] (axis cs:0,2) -- (axis cs:15,2) -- (axis cs:15,0);
\draw[otOrange,line width=0.6pt,dash pattern=on 2pt off 1.5pt] (axis cs:0,3) -- (axis cs:30,3) -- (axis cs:30,0);
\addplot[only marks,mark=*,mark size=1.8pt,otGreenDark] coordinates {(15,2)};
\addplot[only marks,mark=*,mark size=1.8pt,otOrange] coordinates {(30,3)};
\node[pcLegende,anchor=north west,text=otGreenDark] at (axis cs:15.5,0.95) {$t_{1/2}=15$ min};
\node[pcLegende,anchor=north west,text=otOrange!75!black] at (axis cs:30.5,1.75) {$[\mathrm{I_2}]=3{,}0$ mmol$\cdot$L$^{-1}$};
\node[pcLegende,anchor=south east,text=otMuted,fill=white] at (axis cs:90,4.08) {$[\mathrm{I_2}]_\mathrm{f}=4{,}0$ mmol$\cdot$L$^{-1}$};}
\end{axis}
\end{tikzpicture}
